\documentstyle[pstricks,charpath]{article}
%\includeonly{tclcmds}
\title{EcoLab \ Documentation}
\author{Russell K. Standish}
\input{pstricks}
\input{multido}
\input{epsf}
\newcommand{\br}{\mbox{${\bf r}$}} %reproductive rate
\newcommand{\bbeta}{\mbox{\boldmath{$\beta$}}}   %interaction matrix
\newcommand{\bgamma}{\mbox{\boldmath{$\gamma$}}} %dispersal
\newcommand{\bmu}{\mbox{\boldmath{$\mu$}}}       %mutation rate
\newcommand{\bn}{\mbox{\boldmath{$n$}}}          %species density
\newcommand{\nsp}{\mbox{$n_{\rm sp}$}}           %no. of species           

\font\ecolabfont=pagko at 20pt
\newcommand{\EcoLab}{%
\pscharpath*[linewidth=1pt,fillstyle=solid,fillcolor=red]{\ecolabfont
EcoLab}}

\begin{document}
\maketitle
\begin{abstract}
{\sl EcoLab} is a system that implements an abstract ecology
model.\cite{Standish94} It is written as a set of
Tcl/Tk\cite{Ousterhout94} commands so that the model parameters can
easily be changed on the fly by means of editing a script. This
document describes the class libraries that implement the dynamic
array structures, and the interface to the Tcl interpreter, as well
all the Tcl commands and variables that are defined in the {\sl
EcoLab} system. You should also examine the sample {\sl EcoLab} script
(\verb|ecolab.tcl|) provided with the pacakge, and the sample problem
scripts, e.g \verb|pred-prey.tcl| and \verb|lifetime.tcl|.
\end{abstract}

\tableofcontents
\section{The Model}\label{model}

We start with a generalised form of the Lotka-Volterra equation 
\begin{equation}\label{lotka-volterra}
\dot{\bn} = \br*\bn + \bn*\bbeta\bn.
\end{equation}
Here \bn\ is the population density, the component $n_i$ being the
number of individuals of species $i$, \br\ is the difference
between reproduction and death, \bbeta\ is the interaction matrix,
with $\beta_{ij}$ being the interaction between species $i$ and $j$, *
referring to elementwise multiplication and {\tt mutate} is the
mutation operator.

\subsection{Lotka-Volterra Dynamics}

The most obvious thing about equation (\ref{lotka-volterra}) is its
fixed point 
\begin{equation}\label{fixed point}
\hat{\bn} = -\bbeta^{-1}\br,
\end{equation}
where $\dot{\bn}=0$. For this point to be biologically meaningful, all
components of $\hat{\bn}$ must be positive, giving rise to the following
inequalities:
\begin{equation}\label{positive species}
\hat n_i = \left(\bbeta^{-1}\br\right)_i>0, \forall i
\end{equation}
The stability of this point is related to the
negative definiteness of derivative of $\dot{\bn}$ at $\hat{\bn}$. The
components of the derivative are given by
\begin{equation}\label{derivative}
\frac{\partial\dot{n}_i}{\partial n_j} =
\delta_{ij}\left(r_i+\sum_k\beta_{ik}n_k\right) + \beta_{ij}n_i
\end{equation}
Substituting eq (\ref{fixed point}) gives
\begin{equation}
\left.\frac{\partial\dot{n}_i}{\partial n_j}\right|_{\hat{\bn}}=
-\beta_{ij}\left(\bbeta^{-1}\br\right)_i
\end{equation}

Stability of the fixed point requires that this matrix should be
negative definite. Since the $\left(\bbeta^{-1}\br\right)_i$ are
all negative by virtue of (\ref{positive species}), each minor
determinant of this matrix is equal to a minor determinant of \bbeta\
multiplied by a positive number, stability of the equilibrium is
equivalent to \bbeta\ being negative definite.

A weaker condition is to require that the system remain bounded with
time:
\begin{equation}\label{boundedness}
\sum_i\dot{n_i}=\br\cdot\bn + \bn\cdot\bbeta\bn < 0, \;\forall \bn:
\sum_in_i>N \;\exists N
\end{equation}

As \bn\ becomes large in any direction, this functional is dominated
by the quadratic term, so this implies that  $\bn\cdot\bbeta\bn\leq0
\; \forall\bn: n_i>0$. Negative definiteness of \bbeta\ is sufficient,
but not necessary for this condition. For example, the predator-prey
relations (heavily normalised) have the following matrix as \bbeta:
\begin{math}
\bbeta=\left(\begin{array}{cc}
-1 & 2\\
-2 & 0\\
\end{array}\right)
\end{math}
which has eigenvalues $3/2, -5/2$. If we let $\bn=(x,y), x,y\geq0$, then
$\bn\cdot\bbeta\bn=-2x^2$, which is clearly non-positive
for all nonpositive $x$.

Consider adding a new row and column to \bbeta. What is condition is
the new row and column required to satisfy such that equation
(\ref{boundedness}) is satisfied. Break up \bbeta\ in the following
way:
\begin{displaymath}
\left(
  \mbox{
     \begin{tabular}{c|c}
       $\begin{array}{ccc}\ddots\\&{\bf A}\\&&\ddots\end{array}$ & 
       $\begin{array}{c}\vdots\\{\bf B}\\\vdots\end{array}$ \\
       \hline
       $\begin{array}{ccc}\cdots&{\bf C}&\cdots\end{array}$ & D
     \end{tabular}
   }
\right)
\left(
  \mbox{
     \begin{tabular}{c}
     $\begin{array}{c}\vdots\\{\bf n_1}\\\vdots\end{array}$\\
     \hline
     $n_2$
     \end{tabular}
    }
\right)
\end{displaymath}

Condition (\ref{boundedness}) becomes:
\begin{equation}\label{boundedness2}
{\bf n_1}\cdot{\bf A}{\bf n_1} + {\bf n_1}\cdot({\bf B}+{\bf C})n_2 +
Dn_2^2 \leq 0
\end{equation}

Let 
\begin{displaymath}
a=\max_{n=1} \bn\cdot A\bn,\mbox{ and } b=\max_{i}B_i+C_i.
\end{displaymath}
 Then a sufficient but
not necessary condition for condition (\ref{boundedness2}) is
\begin{displaymath}
an_1^2+bn_1n_2+Dn_2^2\leq0
\end{displaymath}

The maximum value with respect to $n_2$ is $an_1^2-(bn_1)^2/4D$, so
this requires that
\begin{equation}\label{boundedness3}
b < 2\sqrt{ad}
\end{equation}

\subsection{Mutation}

With mutation, equation (\ref{lotka-volterra}) reads
\begin{equation}
\dot{\bn} = \br*\bn + \bn*\bbeta\bn + {\tt mutate}(\bmu,\br,\bn).
\end{equation}


The difficulty with adding mutation to this model is how to define the
mapping between genotype space and phenotype space, or in other words,
what defines the {\em embryology}. A few studies, including Ray's
Tierra world, do this with an explicit mapping from the genotype to to
some particular organism property (e.g. interpreted as machine language
instructions, or as weight in a neural net). These organisms then
interact with one another to determine the population dynamics. In
this model, however, we are doing away with the organismal layer, and
so an explicit embryology is impossible. The only possibility left is
to use a statistical model of embryology. The mapping between
genotype space and the population parameters \br,
\bbeta{}  is expected to look like a rugged
landscape, however, if two genotypes are close together (in a Hamming
sense) then one might expect that the phenotypes are likely to be
similar, as would the population parameters. This I call {\em random
embryology with locality}.

  In the simple case of point mutations, the probability $P(x)$ of any
child lying distance $x$ in genotype space from its parent follows a
Poisson distribution. Random embryology with locality implies that the
phenotypic parameters are distributed randomly about the parent
species, with a standard deviation that depends monotonically on the
genotypic displacement. The simplest such model is to distribute the
phenotypic parameters in a Gaussian fashion about the parent's values,
with standard deviation proportional to the genotypic
displacement. This constant of proportionality can be conflated with
the species' intrinsic mutation rate, to give rise another phenotypic
parameter \bmu.  It is assumed that the probability of a mutation generating
a previously existing species is negligible, and can be ignored. We
also need another arbitrary parameter $\rho$, ``species radius'',
which can be understood as the minimum genotypic distance separating
species, conflated with the same constant of proportionality as \bmu.

In summary, the mutation algorithm is as follows:
\begin{enumerate}
\item The number of mutant species arising from species $i$ within a
timestep is $\mu_i\alpha_in_i/\rho$. This number is rounded
stochastically to the nearest integer, e.g. 0.25 is rounded up to 1
25\% of the time and down to 0 75\% of the time.

\item Roll a random number from a Poisson distribution
$e^{-x/\mu+\rho}$ to determine the standard deviation $\sigma$ of phenotypic
variation. 

\item Vary \br\ according to a Gaussian distribution about the
parents' values, with $\sigma\alpha_0$ as the standard deviation,
where $\alpha_0$ is the range of values that \br\ is initialised
to, ie $\alpha_0$=\verb|repro_rate(random,maxval)|$-$\verb|repro_rate(random,minval)|

\item The diagonal part of \bbeta\ must be negative, so vary \bbeta\
according to a log-normal distribution. This means that if the old
value is $\beta$, the new value becomes
$\beta'=-\exp(\log_e(\beta)+\sigma)$. These values cannot become
arbitrarily small, however, as this would imply that some species make
arbitrarily small demands on the environment, and will become infinite
in number. In ecolab, the diagonal interactions terms prevented from
becoming larger than $-r/(.1*{\tt INT\_MAX})$.

\item The offdiagonal components of \bbeta, are varied in a similar
fashion to \br. However new connections are added, or old ones
removed according to $\lfloor 1/r\rfloor$, where $r\in(-2,2)$ is
chosen from a uniform distribution. The values on the new connections
are chosen from the same initial distribution that the offdiagonal
values where originally set with, ie the range
\verb|interaction(offdiag,random,minval)| to
\verb|interaction(offdiag,random,maxval)|. Since $a$ in condition
(\ref{boundedness3}) is computationally expensive, we use a slightly stronger
criterion that is sufficient, computationally tractable yet still
allows ``interesting'' non-definite matrix behaviour namely that the sum
$\beta_{ij}+\beta_{ji}$ should be nonpositive.


\item \bmu\ must be positive, so should evolve according to the
log-normal distribution like the diagonal components of
\bbeta. Similar to \bbeta, it is a catastrophe to allow \bmu\ to
become arbitrarily large. In the real world, mutation normally exists
at some fixed background rate --- species can reduce the level of
mutation by improving their genetic repair algorithms. In ecolab, this
ceiling on \bmu\ is given by the \verb|mutation(random,maxval)| variable.

\end{enumerate}

\section{Structure of the Ecolab Simulation System}

Ecolab is built on three main class libraries: the {\tt arrays} class,
which implements dynamic arrays, the {\tt globals} class, which define
the global state variables of the model and {\tt tcl++} which provides
bindings to the TCL scripting language. Each of these class libraries
will be described in subsequent sections.

Further modules of the system include {\tt analysis} and {\tt
  Xecolab}, which provides a number of generic instruments for
observing the Ecolab system, and {\tt aux}, which contains code for
  implementing  a
client/server version of Ecolab, with the
instruments running in the client system, and the model running on the
server, and code for implementing checkpoint/restart. This minimises the amount of traffic between a compute server
and a client.

The Ecolab model itself is defined in the model specific file {\tt
  ecolab.cc}. To define another model, replace this file by another
  one with similar functions. An example of this is {\tt
  newman.cc}. See \S\ref{new model} for more details.

The whole computation is constructed from a TCL\cite{Ousterhout94}
script. Example scripts include {\tt ecolab.tcl}, {\tt pred-prey.tcl},
{\tt lifetime.tcl} and {\tt console.tcl/engine.tcl} for a sample
client/server system.

\begin{figure}
\input{pst-node}
\pspicture(2,0)(15,7)
\rput(7,7){\rnode{ecolabtcl}{\psframebox{ecolab.tcl}}}
\rput(5,5){\rnode{modeltcl}{\psframebox{model.tcl}}}
\rput(7,5){\rnode{Xecolab}{\psframebox{Xecolab.tcl}}}
\rput(9,5){\rnode{ecolabcc}{\psframebox{ecolab.cc}}}
\rput(5,3){\rnode{analysis}{\psframebox{analysis.cc}}}
\rput(12.5,3){\rnode{globals}{\psframebox{globals.cc}}}
\rput(10,3){\rnode{initarrays}{\psframebox{init\_arrays.cc}}}
\rput(5,1){\rnode{tclpp}{\psframebox{tcl++.h}}}
\rput(7,1){\rnode{arrays}{\psframebox{arrays.h}}}
\rput(7,0){\rnode{carrays}{\psframebox{c\_arrays.c}}}
\ncline{ecolabtcl}{modeltcl}
\ncline{ecolabtcl}{Xecolab}
\ncline{ecolabtcl}{ecolabcc}
\ncline{Xecolab}{analysis}
\ncline{ecolabcc}{tclpp}
\ncline{ecolabcc}{globals}
\ncline{ecolabcc}{arrays}
\ncline{ecolabcc}{initarrays}
\ncline{analysis}{tclpp}
\ncline{globals}{arrays}
\ncline{initarrays}{arrays}
\ncline{arrays}{carrays}
\endpspicture
\caption{Structure of the Ecolab Simulation System}
\end{figure}


\section{Making Ecolab}

Once you have unpacked the gzipped tar file, you should be able to
make ecolab by running make. The distribution will compile ``out of
the box'' on a Linux system --- other systems may require some
modifications to the Makefile in order to get it to work. In
particular, the Makefile supplied assume GNU Make is used. It will
autodetect the operating system (Linux, Solaris and Irix have been
tested), setting options appropriate for that. You may need to adjust
this for local environments, and in particular the Makefile also tests
for machine name, so local variations, such as different include or
library path can be allowed for. In this way, you can use the same
Makefile for sufferent systems. Alternatively, you can set the values
of the variables {\tt FLAGS, LIBS, CC, CPLUSPLUS, LINK} and {ARRAYS},
and comment out the GNU Make logic.

You will need to have TCL/Tk installed. Currently Ecolab has been
developed against TCL 8.0 and Tk 8.0. It can probably be made to work
with later versions with a little effort. You will also need to have
the BLT (version 2.4) in order to use several of the plotting
widgets. Ecolab can be compiled without BLT by commenting out the BLT flags
in the Makefile.

Other options are relatively straight forward. Uncomment the Lapack
flag if you wish to use the {\tt maxeig} command. Uncomment the MPI
flags if you wish to run Ecolab in parallel using the MPI message
passing library.

{\tt CONTIGUOUS} is an optimisation compiler flag that can lead to more
efficient code when the data implementation of arrays are contiguous
(true if running on a single processor workstation; not true if the
arrays are distributed on a distributed memory parallel computer).

\section{Constructing an experiment}

An experiment consists of a TCL\cite{Ousterhout94} script, which binds
the various elements used (the model system, input parameters,
instruments) into an experimental run. An example experiment is {\tt
  ecolab.tcl}. This is an executable script --- once you have made
ecolab, you can run this script.

The script consists of several parts --- the first being a simulation loop
which steps the model through the generate and mutate operators, then
updates the various instruments ({\tt display, plot, connect\_plot}
etc.) A couple of points are worth noting about this: there is a {\tt
  running} flag which controls whether the simulation is running or
not. This is used by the {\bf run} and {\bf stop} buttons to control
the execution of the simulation. 

The other parts of the experimental script have been broken into
separate TCL files --- model.tcl contains the input parameters to the
model, and Xecolab.tcl has the TCL code relating to the X-windows
interface.

\subsection{Input Parameters}\label{input parameters}

The model's parameters are set by TCL variables in {\tt model.tcl}.
The actual data structures of the model are initialised the first time
the model's generate step is called. An example input set is:
\begin{verbatim}
set nsp 10
set dims(0) 2
set density(list) 100
set repro_rate(random) true
set repro_rate(random,seed) 10
set repro_rate(random,maxval) .01
set repro_rate(random,minval) -.005
set interaction(diag,random) true
set interaction(diag,random,seed) 7
set interaction(diag,random,minval) -0.0001
set interaction(diag,random,maxval) -0.00005
set interaction(offdiag,random) true
set interaction(offdiag,random,minval) -0.0001
set interaction(offdiag,random,maxval) 0.0001
set interaction(offdiag,random,connectivity) 2
set interaction(offdiag,random,seed) 12
set mutation(random) true
set mutation(random,maxval) .5
\end{verbatim}
The initial dimension of the system is 10, set by the variable {\tt
  nsp} for number of species. The spatial dimensions are set by
assigning to the array variable dims --- in this case, the system is
one dimensional, with 2 cells (see \S\ref{spatial}). The
remaining variables refer to the initial values of the array variables
in Ecolab. ${\tt density}\equiv\bn$ is assigned the value 100 to each
element. One can assign a TCL list to these variables also --- see
{\tt pred-prey.tcl} for an example. The other variables ${\tt
  repro\_rate}\equiv\br$, ${\tt interaction}\equiv\bbeta$ and ${\tt
  mutation}\equiv\bmu$ are assigned uniformly random values in this
case. For a full rundown on the syntax of the input parameters, see
section \ref{init_arrays}.

\subsection{The main button bar}

The main button bar has three predefined buttons, {\bf quit} (which is
bound to the TCL command {\tt exit\_ecolab} for leaving the
application), {\bf run}, which calls the {\tt simulate} procedure, and
{\bf stop} which sets the {\tt running} flag to zero, suspending the
experiment. Additionally, there are three user defined buttons, which
can have functions bound to them by adding TCL code after {\tt source
  Xecolab.tcl} like  the following example:
\begin{verbatim}
set more_ecolab 1
source Xecolab.tcl
.user1 configure -text condense -command condense
map_mainwin
\end{verbatim}
The \verb+more_ecolab+ variable prevents Xecolab from mapping the
mainwindow, allowing further configuration information to be
provided. It is then the user responsibility to map the window using
\verb+map_mainwin+. 

\subsection{Instrumentation}

Currently four instrumented widgets are included in the Ecolab
distribution. The first time these widgets are called, they
instantiate themselves in a separate window. This follows the
philosophy that initialisation should be a transparent operation. A
number of these widgets depend on the BLT toolkit, so it is wise to
include BLT at the compilation stage. Each widget is independent, and
so multiple instruments can be operating at the same time.

\subsubsection{Plot}

Plot a number of TCL variables

Usage:

{\tt plot} {\em plotname} [{\tt -title} {\em title string}] {\em x y1}
[{\em y2\ldots}]

Note that the arguments are actual TCL variable names, not values
(i.e. don't use \${\em variable} syntax). {\em plotname} labels the
particular widget, so multiple plots can be performed simultaneously.

The plot can be zoomed by selecting a rectangular region using the
left mouse button. The right mouse button reverts to the original scale.

\subsubsection{Histogram}

Usage:

{\tt histogram}  {\em plotname} [{\tt -title} {\em title string}] {\em
  list of values}

This command adds the values to a histogram plot. A record of all data
values so far added to the histogram is stored in the file {\em
  plotname}{\tt .dat}. This allows automatic recalculation of the
histogram's bins whenever a data value goes out of bounds. The number
of bins is also dynamically controlled by the scale widget on the
side. The histogram is recalculated next time {\tt histogram} is
called after the number of bins has been altered. This can be a time
consuming business, so it is recommended that the experiment is
suspended before changing the number of bins.

There is also the option of plotting the histogram with logarithmic scales.

\subsubsection{Display}\label{display}

Usage:

{\tt display} {\em cell} 

This command displays species density in the {\em cell} as a function
of time. The colour of the lines used are controlled by the {\tt
  palette} TCL variable.(\S\ref{palette}) This is a list of X-window
colours to use.

\subsubsection{Connect\_plot}\label{connect}

This command takes no arguments, and displays the connectivity of the
interaction matrix for each cell. The ecologies are ordered within
each cell to show up the independent subecologies, and if the {\tt
  palette} variable (\S\ref{palette}) has been defined, each of the independent
subecology is coloured differently.

There are two buttons which allow the user to zoom in and out by a
scale factor of two.  One can also zoom in by clicking with the right
mouse button at the location you wish to zoom in on.  The plot can
also be dragged with the left mouse button to change the view.

\subsection{Auxilliary Commands}

\subsubsection{get\_global\_vars/data\_server}

Syntax:

{\tt get\_global\_vars} {\em server port}

{\tt data\_server} {\em port}

This pair implements a client server connection. {\tt data\_server} is
executed on the compute server, and services any requests coming in on
the given port. {\tt get\_global\_vars} attaches to the compute
server, and downloads the compute servers global state variable into
the clients global variables. The client can the follow on with the
usual instruments for analysing the data. The advantage in this
approach is that  X-window traffic can be avoided, the total amount of
traffic between client and server can be controlled by how often these
routines are called. An example setup is located in {\tt console.tcl}
and {\tt engine.tcl}. An alternative client/server scenario can be
contructed using Tcl sockets, and the console2.tcl/engine.tcl give an
example of just transferring the \bn{} of the predator-prey
example. This has a great deal of flexibility, allowing, for example,
messages to be propagated from client back to the server to allow user
interactivity into the model

\subsubsection{checkpoint/restart}

Syntax:

{\tt checkpoint} {\em filename}

{\tt restart} {\em filename}

These commands dump the contents of the global variables into a file,
and correspondingly reload the global state variables from the file,
in order to implement checkpoint-restart for a batch, or long running
environment. 

\subsection{Ecolab Model commands}

\subsubsection{generate}

This implements the basic Lotka-Volterra equations:

\begin{math}
\dot{\bn} = \br*\bn + \bn*\bbeta\bn
\end{math}

with \br\ being the reproduction rate and \bbeta\ being the
interspecies interaction. This is implemented as a single line:

\begin{verbatim}
  density += repro_rate * density + (interaction * density) * density;
\end{verbatim}

This command also increments the timestep counter {\tt tstep}.

\subsubsection{condense}\label{condense}

This command compacts the systems of equations by removing extinct
species where $n_i=0$.

\subsubsection{mutate}\label{mutate}

This applies the point mutations to the system. In this case, the
number of mutants each species produces is calculated from
(\verb|sp_sep|)\br\bmu\bn. For each of these new species, mutate
calculates a ``genetic drift'' probabilistically according a Poisson
distribution with mean given by \bmu. It then arranges for the values
of \br, \bbeta\ and \bmu\ to be mutated according to a Gaussian
distribution, with standard error given by the ``genetic drift''. The
skeleton of the code follows:

\begin{verbatim}
  /* calculate the number of mutants each species produces */
  new_sp = (double) sp_sep * repro_rate * mutation * density;

  /* generate index list of old species that mutate to the new */
  new_sp = gen_index(new_sp); 

  /* collect the old phenotypes */
  new_repro_rate = repro_rate[new_sp];
  new_mutation = mutation[new_sp];
... and the same for interaction, but messy ...

  /* calculate the genetic distances for the mutants from the parents */
  array gdist(new_sp.size);
  gdist.fillprand();
  gdist *= new_mutation;

  /* 
    change phenotypes randomly, according a normal distribution 
    with std dev gdist 
  */
  gspread( new_repro_rate, gdist );
  gspread( new_interaction, gdist );
  gspread( new_mutation, gdist );

  /* concatonate new species */
  density <<= new_sp;
  repro_rate <<= new_repro_rate;
  interaction <<= new_interaction;
  mutation <<= new_mutation;

  nsp += new_sp.size;
\end{verbatim}

The offdiagonal components are varied in a similar way, except that there
is also a finite probability that new connections are added or removed in
the mutant species.

\subsubsection{migrate}\label{migrate}

This operator implements migration within cellular Ecolab. This
updates density values according to the difference with the 4 nearest
neighbours: $\bn+=0.25*(\bn_n+\bn_e+\bn_s+\bn_w)-\bn$, where the
$n,e,s,w$ index the north, east, south and west neigbouring cells.

\subsubsection{maxeig}\label{maxeig}

\verb|maxeig| returns the maximum eigenvalue of \bbeta. If this number
is negative, the equilibrium point is stable, if positive, it is
unstable. As reported in Standish (1994)\cite{Standish94}, the mutation
drives the maximum eigenvalue slightly positive, then instabilities
act to push the eigenvalue back to zero. This routine is not multicellular.

\subsubsection{lifetimes}\label{lifetimes}

\verb|lifetimes| records the timestep when a
species passes a threshold (hardwired at 10) in the {\tt create}
iarray. If a species has yet to pass the threshold, or has gone
extinct, the value in {\tt create} is zero. Upon return, this routine
returns the lifetime of the species that have gone extinct. This can
then be passed to a histogram routine, or written to a file.

\subsubsection{get}\label{TCLget}

Syntax:

{\tt get} {\em variable}

where {\em variable} is the name of an ecolab state variable from the
list: nsp, tstep, density, tags, species, repro\_rate, mutation,
interaction(diag), interaction(val), interaction(row),
interaction(col). This command returns the value of the state variable
a TCL list, which can be further processed by TCL programs.


\section{Creating New Instruments}

The best way to proceed is to copy an existing instrument, and modify it
to your needs. The interface to the TCL language can be done almost
entirely through the {\tt tcl++} class library (\S\ref{tcl++}). Use the
{\tt NEWCMD} macro to instantiate a new TCL command. This may be the
complete widget, or merely a component of it. One of the first operations
should be to create a new window through the {\tt
  newwin}(\S\ref{TCLnewwin}) TCL command. After that, the instrument may
be encoded entirely in TCL code, using the {\tt get}(\S\ref{TCLget}) call
to return state variables as TCL lists. Alternatively, you may write C++
code to access the state variables directly through the array class. The
{\tt ecolab.h} header defines references to the generic state variables
as:
\begin{verbatim}
extern global global_vars;
extern int &ocell, &ncells, &tstep;
extern iarray &nsp, &dims;
\end{verbatim}

Furthermore, the following model specific variables are defined in {\tt
  ecolab.cc} for the ecolab model outlined in \S\ref{model}.
\begin{verbatim}
extern iarray &density;
extern iarray &create;
extern iarray &species;
extern array &repro_rate;
extern array &mutation;
extern sparse_mat &interaction;
\end{verbatim}
The references themselves point to the appropriate members of {\tt
  global\_vars}. In terms of the model defined in \S\ref{model}, {\tt
  ocell} refers to the origin cell on this processor (this is zero on
processor 0 e.g. on a uniprocessor). {\tt ncells} refers to the number
of cells on this processor, and {\tt tstep} the timestep (which
incremented on each call to generate). {\tt nsp} is an array (of size
{\tt ncells}), that contains the list of \nsp{}s for the cells on the
processor. {\tt dims} contains the array of dimension sizes for
mapping the cell structure back to physical space. see
\S\ref{spatial}. {\tt density} is \bn, {\tt repro\_rate} is \br, {\tt
  mutation} is \bmu{} and {\tt interaction} is \bbeta. {\tt species}
is used to enumerate species (position within the arrays don't count,
as after condensing, extinct species are removed from the system).
{\tt create} is used in the {\tt lifetimes}(\S\ref{lifetimes})
command.

\section{Creating a New Model}\label{new model}

The code that explicitly defines the ecolab model is contained in {\tt
  ecolab.cc}. An example of a completely different model is contained in
{\tt newman.cc}, and {\tt newman.tcl}.  This implements a generalisation
of a model described by \cite{Newman97} where the species density is a
boolean value (either a species is extinct, (value=0) or it is alive). The
dynamics are controlled by a set of threshhold values, which are different
for different species, and a randomly generated stress $\eta$ which
follows some distribution $P_{\rm stress}(\eta)$ (eg Gaussian). If the
stress exceeds that of the species' threshold, then the species becomes
extinct. There is a small additional probability that a species becomes
extinct even when the threshold is not extinct (decreasing proportionally
to the amount that the threshold is greater than the stress). This
prevents the system becoming stuck with maximally fit organisms (maximal
thresholds). Mutation is performed in a similar way to ecolab, although
there is no need to consider an interaction matrix. The number of mutants
is computed such that the average no. of mutations equals the average
number of extinctions. This is to keep the number of species roughly
constant in time, although it can fluctuate around the mean. This is a
generalization of Mark Newman's constant species condition, and appears to
be necessary to avoid the system exploding or crashing to zero.

The main items that need to be constructed in a model file are:
\begin{enumerate}
\item declaration of the state variables
\item A subroutine to initialize the state variables {\tt
    init\_global\_vars}, that call routines in {\tt init\_arrays.cc}\S\ref{init_arrays}.
\item A TCL command to step the simulation, and perform mutations.
\item A TCL command for housekeeping roles, such as {\tt condense}
\item TCL interfaces to widgets such as {\tt display} or {\tt connect\_plot}
\item TCL interface for returning state variables as lists {\tt get}
\item Any other more explicit computations, such as {\tt lifetimes}
  or {\tt max\_eig}.
\end{enumerate}


\section{Spatial Variation}\label{spatial}

The ecolab model can run in multicellular mode by assigning to the
{\tt dims}. The model in equation \ref{lotka-volterra} is replicated
in each cell. The idea behind this is to set up a spatially dependent
model. In case of two dimensional space, one can specify the x and y
dimensions by assigning the following values as part of the input
parameters \S\ref{input parameters}:
\begin{verbatim}
set dims(x) 10
set dims(y) 20
\end{verbatim}
This specifies a $10\times20$ grid of 200 cells.
The array indices ``x'' and ``y'' are purely arbitrary, one could use
``0'' and ``1'' instead.

In the multicellular case, the number of species in each cell is set
to be the value of the TCL variable {\tt nsp}. In order to specify
different numbers of species in each cell at initialisation, one can
assign to the {\tt nsp} array, in the following syntax (eg a 2D
system):
\begin{verbatim}
set nsp(0,0) 10
set nsp(0,1) 12
set nsp(1,0) 15
set nsp(1,1) 13
\end{verbatim}

One can also initialise other state variables with spatial dependence
by appending cell indices, e.g.
\begin{verbatim}
set density(list,0,0) 100
set density(list,0,1) 100
set density(list,1,0) 100
set density(list,1,1) 100
\end{verbatim}
In this case, one may need to write a short TCL loop to intialise a
particular pattern. This feature is not yet implemented as of version 3.2

Migration is performed using the {\tt migrate}\ref{migrate} operator.

\subsection{Parallel Execution}

A multicellular model can be run on an MPI parallel system by
compiling with the MPI flags turned on. Ecolab distributes the cells
across the processors in a block fashion. A TCL command can be
written that executes on all processors by placing the {\tt PARALLEL}
directive at the start of the command. Otherwise (by default) all code
executes just on processor 0. An example of a useful utility function is
{\tt get\_all\_globals} which assembles a complete representation of
the whole system on processor 0. This uses the {\tt PARALLEL give\_me}
command to obtain the slave processor state variables. There is a
corresponding {\tt take\_it} command to push state variable out to all
the slaves. These expect the master process to broadcast the data, so
these routines should not be called from TCL scripts, otherwise Ecolab
will hang.

Work on using the spatial Ecolab is reported in
\cite{Standish98c}. The model showed good scalability, and even
superlinear speedup in some circumstances.

The parallel coding practices employed in Ecolab are experimental, and
may well change in the next release.

\section{tcl++}\label{tcl++}

This section describes a light weight class library for creating
Tcl/Tk applications. \cite{Ousterhout94} It consists of a header file
\verb|tcl++.h|, and a main program \verb|tclmain.cc|. All you need to
do to create a Tcl/Tk application is write any application specific
Tcl commands using the NEWCMD macro, link your code with the
\verb|tclmain| code, and then proceed to write the rest of your
application as a Tcl script. The \verb|tclmain| code uses argv[1] to
get the name of a script to execute, so assuming that the executable
image part of your application is called \verb|appl.exe|, you would
start off your application script with
\begin{verbatim}
#!appl.exe
...Your Tcl/Tk code goes here...
\end{verbatim}
then you script becomes the application, directly launchable from the
shell.

\subsection{Global Variables}

The following global symbols are defined by \verb|tcl++|:
\begin{description}
\item[interp] The default interpreter used by Tcl/Tk
\item[mainWin] The window used by Tk, of type \verb|Tk_Window|, useful
for obtaining X information such as colourmaps
\item[TclError] A longjmp label used for trapping errors --- this will
be replaced by exception handling once g++ has this ability
\item[error] To flag an error, call \verb|error(char *format,...)|.
The format string is passed to sprintf, along with all remaining
arguments. The error message is attached to Tcl's error log, and then
control branches to \verb|TclError|, which returns a \verb|TCL_ERROR|
code back to Tcl.
\end{description}

\subsection{Creating a Tcl command}
To register a procedure, use the \verb|NEWCMD| macro. It takes two
arguments, the first is the name of the Tcl command to be created, the
second is the number of arguments that the Tcl command takes.
\verb|NEWCMD| does the following:
\begin{itemize}
\item lodges  the commands with the interpreter prior to the
execution of tclmain,
\item sets up the longjmp label to trap errors,
\item checks that the declared number of arguments match the actual
number (type checking is beyond the ANSI C Preprocessor!)
\item executes the contents of the macro \verb|GLOBAL_INIT_HOOK|,
which can be set to (for example) initialize any global variables from
Tcl variables.
\item declares a void returning C++ function of the same name as the
Tcl command. The arguments are passed via the argc, argv
mechanism. All that is required is the body of the function to be put
in.
\end{itemize}

Here is an example.
\begin{verbatim}
#include "tcl++.h"

#undef GLOBAL_INIT_HOOK
#define GLOBAL_INIT_HOOK init_funny();

static int funny;

void init_funny()  /* this must be declared before NEWCMD */
{
   if ( 
     Tcl_GetInt(interp, 
        Tcl_GetVar(interp,"funny",TCL_GLOBAL_ONLY), 
        &funny) 
     !=TCL_OK) error("funny not found"); 
}


NEWCMD(silly_cmd,1)
{printf("argv[1]=%s funny=%d\n", argv[1], funny);}
\end{verbatim}
This creates a Tcl command that prints out its argument, and the value
of the variable \verb|funny|.

\subsection{tclvar}

Tcl variables can be accessed by means of the tclvar class. For
example, if the programmer declares:

\verb|tclvar hello="hello"; float floatvar;|

then the variable hello can be used just like a normal C variable in
expression such as 

\verb|floatvar=hello*3.4;|.

The class allows assignment to and from \verb|double| and \verb|char*|
variables, incrementing and decrementing the variables, compound
assignment and accessing array Tcl variables by means of the C array
syntax. The is also a function that tests for the existence of a Tcl
variable.

\pagebreak[4]
The complete class definition is given by:

\begin{verbatim}
class tclvar
{ 
 public:

/* constructors */
  tclvar();
  tclvar(char *nm, char* val=NULL);
  tclvar(tclvar&);
  ~tclvar();

/* These four statements allow tclvars to be freely mixed with arithmetic 
 expressions */
  double operator=(double x);
  char* operator=(char* x);
  operator double ();
  operator char* ();

  double operator++();
  double operator++(int);
  double operator--();
  double operator--(int); 
  double operator+=(double x);
  double operator-=(double x);
  double operator*=(double x);
  double operator/=(double x);

  tclvar operator=(tclvar x); 
/* arrays can be indexed either by integers, or by strings */
  tclvar operator[](int index);
  tclvar operator[](char* index);

  friend int exists(tclvar x);
};
\end{verbatim}

\subsection{tclcmd}

a tclcmd allows Tcl commands to be executed by a simple \verb|x <<|
{\em Tcl command} syntax. The commands can be stacked, or
accumulated. Each time a linefeed is obtained, the command is
evaluated.

\pagebreak[4]
For example:
\begin{verbatim}
tclcmd cmd;

cmd << "set a 1\n puts stdout $a";
cmd << "for {set i 0} {$i < 10} {incr i}"
cmd << "{set a [expr $a+$i]; puts stdout $a}\n";
\end{verbatim}

The behaviour of this command is slightly different from the usual
stream classes, as a space is automatically inserted between arguments
of a \verb|<<|. This means that arguments can be easily given - eg
\begin{verbatim}
cmd << "plot" << 1.2 << 3.4 << "\n";
\end{verbatim}
will perform the command {\tt plot 1.2 3.4}. In order to build up a
symbol from multiple arguments, use the append operator. So
\begin{verbatim}
cmd << "plot"|3 << "\n";
\end{verbatim} 
executes the command {\tt plot3}.

The full class is given by:
\begin{verbatim}
class tclcmd
{
 public:
  char *result; /* The result of the previous command is placed here */
  tclcmd();
  ~tclcmd();
  tclcmd& operator<<(char* cmd);
  tclcmd& operator<<(double x);
  tclcmd& operator<<(int x); 
  tclcmd& operator<<(unsigned x);
  tclcmd& operator<<(tclvar x);
  tclcmd& operator|(char* cmd);
  tclcmd& operator|(double x);
  tclcmd& operator|(int x); 
  tclcmd& operator|(unsigned x);
  tclcmd& operator|(tclvar x);
};
\end{verbatim}

\subsection{tclreturn}

The {\tt tclreturn} class is used to supply a return value to a TCL
command. It inherits from {\tt ostrstream}, so has exactly the same
behaviour as that class. When this variable goes out of scope, the
contents of the stream buffer is written as the TCL return value. It is
an error to have more than one {\tt tclreturn} in a TCL command, the
result is undefined in this case.

\subsection{tclindex}

This class is used to index through a TCL array. TCL arrays can have
arbitrary strings for indices, and are not necessarily ordered. The
class has three main methods: {\tt start}, which takes a {\tt tclvar}
variable that refers to the array, and returns the first element in
the array; {\tt incr}, which returns the next element of the array;
and {\tt last} which returns 1 if the last element has been read. A
sample usage might be (for computing the product of all elements in
the array {\tt dims}):
\begin{verbatim}
tclvar dims="dims";
tclindex idx;
for (ncells *= (int)idx.start(dims); !idx.last(); ncells *= (int)idx.incr() );
\end{verbatim}

The class definition is given by:
\begin{verbatim}
class tclindex  
{
public:
  tclindex();
  tclindex(tclindex&);
  ~tclindex() {done();}
  tclvar start(tclvar&);
  inline tclvar incr();
  tclvar incr(tclvar& x) {incr();}  
  int last();
};
\end{verbatim}






\section{arrays}

This section documents the array and related data types. These are
dynamically sized, with the size stored at all times in the size
member. These classes are written as a C++ wrapper around a C Library
{\tt c\_arrays}. This enables the {\tt c\_arrays} library to be
implemented to make use of parallel computation and an example
implementation in C* ({\tt cs\_arrays.cs}) is provided. Unfortunately,
the function \verb|offmul| is actually the computationally dominant
routine, and this proves to be the most difficult to optimise for
parallel computers.

\subsection{array and iarray}\label{array}

\begin{verbatim}

class array
{
 public:

  int size; 
  void *list;  

  array (int s=0);
  array ( array & x);
  ~ array ();
  array  operator=( array  x); 

  operator double *();
  array  operator=( double  * x);

  /* unary minus */
  array  operator-();

  /* binary arithmetic operators */
  array   operator  +  (  array   x);
  array   operator  -  (  array   x);
  array   operator  *  (  array   x);
  array   operator  /  (  array   x);

  /* binary comparison operators */
  iarray  operator  <  (  array   x);
  iarray  operator  <=  (  array   x);
  iarray  operator  >  (  array   x);
  iarray  operator  >=  (  array   x);
  iarray  operator  ==  (  array   x);
  iarray  operator  !=  (  array   x);

  /* concatenation */
  array  operator << ( array  x);

  /* broadcast arithmetic ops - commutative equivalents defined
  elsewhere */
  array   operator  +  (  double   x);
  array   operator  -  (  double   x);
  array   operator  *  (  double   x);
  array   operator  /  (  double   x);

  /* broadcast comparison ops  - commutative equivalents defined
  elsewhere */
  iarray  operator  <  (  double   x); 
  iarray  operator  <=  (  double   x);
  iarray  operator  >  (  double   x);
  iarray  operator  >=  (  double   x);
  iarray  operator  ==  (  double   x);
  iarray  operator  !=  (  double   x);

  /* concatenate an element  - commutative equivalents defined
  elsewhere */
  array  operator << ( double  x);

  /* compound assignment */
  array   operator  += (  array   x);
  array   operator  += (  double   x); /* broadcast add to */
  array   operator  -= (  array   x);
  array   operator  -= (  double   x);
  array   operator  *= (  array   x);
  array   operator  *= (  double   x);
  array   operator  /= (  array   x);
  array   operator  /= (  double   x);
  array   operator  <<= (  array   x); /* compound concatenation */ 
  array   operator  <<= (  double   x);

  /* index operations */
#ifdef CONTIGUOUS
  double&   operator[](int);
#else
  par_addr_double operator[](int);
#endif
  par_addr_array  operator[](iarray); /* gather/scatter ops */

/* all of the previous routines have been defined for the iarray class
as well, with every occurance of array being replace by iarray, and
every occurance of double being replaced by int */

/* multiply an array by an iarray */
  inline array operator*(iarray);

/* broadcast an integer to an array */
  array operator=(int x);

/* probabilistically round an array to an iarray. See rounding below */
  operator iarray();
};  

class iarray
{
 public:
  /* all the previous methods defined for array ar also defined
     equivalently for iarray */

  /* extra iarray methods */
  /* boolean operators */
  iarray operator ! ();
  iarray operator  && (iarray x);
  iarray operator || (iarray x);
  iarray operator&=(iarray x) {return *this = *this && x;}
  iarray operator|=(iarray x) {return *this = *this || x;}

 /* probabilistic assignment - see rounding below */
  void operator=(array x);    
  void operator+=(array x);  

 /* convert to array */
  inline operator array();
}

\end{verbatim}

The implementation of the data storage is pointed to by
\verb|list|. More than one array variable may have the same
\verb|list| pointer, so that assignment does not involve copying of
the data, just copying the list pointer. This allows arrays to be
passed by value efficiently in the C++ code. In the case where {\tt
  list} points to a contiguous memory location where the data is
stored, defining {\tt CONTIGUOUS} allows obvious optimisations for
index operations.

iarray is the integer version of array. Binary operations +,*,=
etc. have been defined elementwise, otherwise a scalar value may be
broadcast over the whole array. The concatenation operation joins two
arrays to make a larger array, e.g. $\{1,2,3\} << \{4,5,6\} =
\{1,2,3,4,5,6\}$. 

\subsubsection{Indexing and Gather/Scatter}

The expression \verb|a[10]| is meant to be an {\em lvalue} (i.e. can
be assigned to) that refers to the tenth element of a. When the values
of a are internally stored in a contiguous memory location, this is
easy, with the method returning a reference. However, if the data is
stored in a distributed fashion, references are not possible, and a
function will need to be called to extract the value, or alternatively
alter the value, depending on the context. This achieved by a helper
class \verb|par_addr_double|. This can either be type cast to a
\verb|double|, in which case the function returning the value of
\verb|a[i]| is called, alternatively, it can be assigned to, in which
case the alteration function is called. If more than one array
variable refers to the same underlying data, a copy is made of the
data before an array element is altered.

A {\em gather} operation, or {\em vector index} operation is expressed
by \verb|a[i]| where {\tt i} is an {\tt iarray} variable. The
expression \verb|a[i]=x| {\em scatters} the values of {\tt x} into the
array {\tt a} according to the vector {\tt i}. Naturally, the length
of {\tt i} and {\tt x} must agree. A fresh copy of {\tt a} data is
made if needed.


\subsubsection{Rounding}

The assignment of an array variable to an iarray variable is
interesting, in that the results are rounded up or down
probabilistically. If a value 0.2 is to be assigned, then it has an
80\% chance of being rounded down to 0, and a 20\% chance of being
rounded up to 1.

\subsection{{\tt sparse\_mat}}\label{interact_array}

The {\tt sparse\_mat} data type is designed to hold the
interaction matrix \bbeta, which is generally a sparse matrix. It is
built up of array and iarray components:

\begin{verbatim}
class interact_array
{
 public:
  array diag, val;
  iarray row, col;
  sparse_mat(int s=0, int o=0)
    {diag=array(s); val=array(o); row=iarray(o); col=iarray(o);}
  array operator*(iarray& x);  /* matrix multiplication */
  sparse_mat operator=(sparse_mat x);
};
\end{verbatim}

{\tt diag} stores the diagonal components of the array, {\tt val} is
the packed list of offdiagonal values, with {\tt row} and {\tt col}
being the index lists. The only important operator defined for this
class is the matrix operation, which is defined as
\begin{verbatim}
beta*x == beta.diag*x + (beta.val*x[beta.col])[beta.row]
\end{verbatim}
but is implemented separately for efficiency reasons.

It is possible to represent the offdiagonal array differently  for
efficiency reasons. For example, if it is desired to represent the
offdiagonal elements as a dense 2D array, one can create an extra
pointer in the underlying implementation of array. Most of the time,
this pointer is NULL, but when the \verb|cs_arrays| routine
\verb|offmul| is called, it will check this pointer for the {\tt val}
array. If it is NULL, it will create the efficient
representation, otherwise it will reuse the existing one. Because the
only way the array's actual value will get out of synch with the
efficient representation is by index assignment, \verb|put_double()|
and \verb|put_double_array| (as well as obviously
\verb|delete_array()|) will need to be modified to deallocate the
efficient representation.

\subsection{Global functions}

In all of the following, if it reasonable for an iarray version to
exist, then it will.

\begin{itemize}
\item array {\tt merge}( iarray {\em mask}, array {\em a}, array {\em b} )

return an array (or iarray) {\em r} such that where {\em mask[i]==1},
{\em r[i]=a[i]}, otherwise {\em r[i]=b[i]}.

\item array {\tt pack}( array {\em x}, iarray {\em mask}, [int {\em
ntrue}])

Construct a new array from elements of  {\em x} that correspond to
where the mask is true. {\em ntrue} is an optional parameter equal to
the number of true elements of {\em mask}. If not given, then {\tt
pack} will count the true elements of {\em mask}. Give this parameter
if you are using the same {\em mask} in multiple {\tt pack} statements.

\item iarray {\tt enumerate}(iarray {\em mask})

Return the running sum of mask.

\item iarray {\tt pcoord}( int {\em size})

return [0..size-1]

\item iarray {\tt gen\_index}(iarray {\em x}) 

generate a list of species numbers, with each number appearing in the
list according to the value passed in its position. For example, if
{\em x}=\{0,0,1,2,0,1\} then the output from this program will be
\{3,4,4,6\}


\end{itemize}

\subsubsection{Reduction functions}

\begin{itemize}
\item double {\tt max}( array {\em x}, [iarray {\em mask}]) 
\item double {\tt min}( array {\em x}, [iarray {\em mask}]) 
\item double {\tt sum}( array {\em x}, [iarray {\em mask}]) 

Return the maximum, minimum and sum of an array, optionally masked.

\end{itemize}

\subsubsection{Array functions}
\begin{itemize}
\item array {\tt abs}(array {\em x})
\item iarray {\tt sign}(array {\em x})
\item array {\tt exp}(array {\em x})
\item array {\tt log}(array {\em x})
\end{itemize}

\subsubsection{Random number functions}

\begin{itemize}
\item void {\tt fillrand}(array {\em x}) 

Fill {\em x} with random numbers from the uniform distribution over $[0,1]$

\item void {\tt fillprand}(array {\em x})

Fill {\em x} with random numbers from the Poisson distribution $e^{-x}$

\item void {\tt fillgrand}(array {\em x})

Fill {\em x} with random numbers from the Gaussian distribution $e^{-x^2/4}$

\end{itemize}

\subsubsection{Stream Functions}
\begin{itemize}
\item \verb|ostream& operator<<(ostream& s, array x)|
\item \verb|ostream& operator<<(ostream& s, iarray x)|
\item \verb|ostream& operator<<(ostream& s, par_addr_double x)|
\item \verb|ostream& operator<<(ostream& s, par_addr_int x)|
\end{itemize}


\include{init_arrays}
\section{Globals class}

The {\tt globals} class is used to refer to all state variables as one
object, for the puposes of things like checkpoint restart. The object
may be a subrange of cells, for example the cells that are represented
on a particular processor. Methods exist to extract or replace
subranges of cells. The index operator allows components of the arrays
to be extracted symbolically. So if {\tt density}$\equiv${\tt
  global\_vars.iarrays[0]}, then {\tt g[density]} returns {\tt g.iarrays[0]}.

\begin{verbatim}
//A class for storing the global variables used in the model
class global
{
public:
  //origin cell and number of cells on this processor 
  int ocell, ncells, tstep; 

  // dimensions array and number of species in each cell
  iarray dims, nsp, *coords;  

  // list of iarrays
  iarray *iarrays;

  //list of arrays
  array *arrays;

  //list of sparse_mats
  sparse_mat *sparse_mats;

  //lengths of above lists
  int niarr, narr, nspars;

  global(); 
  /* i=no. of iarrays, a=no. of arrays and s=no. of sparse_mats */
  global(int i, int a, int s);
  global(global&); 
  global operator=(global);
  ~global();

  /* get local equivalents of a global variable */
  iarray& operator[](iarray& var); 
  array& operator[](array& var);
  sparse_mat& operator[](sparse_mat& var);
  

  /* return just the global variables belonging to cell, or a number
      of cells cell, cell+1, ... cell+nc-1 */
  global get_cell_vars(int cell, int nc=1);
 
  //replace cells variables by those contained in vars
  void put_cell_vars(global vars);
  //add a cell into this global
  void append(global vars);

  //pack up the global data into a contiguous package
  void packup(glue& buffer);
  //unpack contiguous representation of global data
  void unpack(glue& buffer);
};

\end{verbatim}


The {\tt glue} class is a helper class that represents the global
variables contiguously. It has four important members; {\tt char
  *data} and {\tt int size}, which point to the contiguous
representation, and the number of bytes within that representation;
and the methods \verb|<<| and \verb|>>|, which appends/extracts an
object into/from the representation. The {\tt packup} and {\tt unpack}
methods within the globals class convert between the glue
representation and the globals representation.

\begin{verbatim}
//A Class for grouping a whole bunch of binary things together
class glue
{
public:
  char *data;
  int size;
  glue(int sz);
  ~glue()
  void operator<<(int);
  void operator<<(double);
  void operator<<(iarray);
  void operator<<(array);
  void operator<<(sparse_mat);
  void operator>>(int&);
  void operator>>(double&);
  void operator>>(iarray&);
  void operator>>(array&);
  void operator>>(sparse_mat&);
};

\begin{verbatim}
//Assemble a full set of global variables by combining data on each processor
void get_all_globals(global&);
\end{verbatim}

collect the global variables on each processor onto processor 0.

\section{Analysis}\label{analysis}

\subsection{newwin}\label{TCLnewwin}

Syntax:

{\tt newwin} {\em name}

Creates a new toplevel window for a TCL widget, with Tk identifier
{\em name}. {\em name} must start with a full stop.

\subsection{Reduction Functions}

Syntax:

{\tt max} {\em tcllist}\\
{\tt min} {\em tcllist}\\
{\tt av} {\em tcllist}\\

The maximum, minimum and average of a TCL list.

\subsection{Palette Variable}\label{palette}

This variable is used by both the {\tt display\_stub} and {\tt
  connect\_stub} routines to define a palette of colours for colouring
  the species in the instruments. If the TCL variable {\tt palette} is
  assigned a list of X-windows colours (on many systems, a list of
  such colours is found in {\tt /usr/lib/X11/rgb.txt}),
  then the palette class can be used like an array within C++,
  returning the colour name as a string:
\begin{verbatim}
palette[i]
\end{verbatim}
returns the \verb|i%n|th colour, where {\tt n} is the number of
colours in the palette list.

\subsection{display\_stub}

Syntax:

\begin{verbatim}
void display_stub(iarray& density, iarray& species, iarray& nsp, 
                  int ncells=1, int cell=0);
\end{verbatim}

This function implements the {\tt display}(\S\ref{display}) widget.
The value of each component of density is plotted as a function of
{\tt tstep}, and coloured by the index {\tt species}. Set {\tt
  ncells>1} and {\tt cell} in order to display only a single cell, as
in ecolab. These variables may be ignored if not meaningful. {\tt nsp}
must be set with the cellular structure used by {\tt density}. see
\S\ref{spatial}. For a non-cellular model, it contains just the 1
component, equal the {\tt density.size}.

\subsection{connect\_stub}

Syntax:

\begin{verbatim}
void connect_stub(iarray& density, sparse_mat& interaction, iarray& nsp, 
                  int ncells=1)
\end{verbatim}

This function implements the {\tt connect\_plot}(\S\ref{connect})
widget. The {\tt row} and {\tt col} connections are plotted, and
sorted and coloured according to connected subsets. {\tt density} is
only used to indicate extinctions (where {\tt density[i]==0}, the row
and column is coloured in wheat. Set {\tt ncells>1} and {\tt cell} in
order to display only a single cell, as in ecolab. These variables may
be ignored if not meaningful. {\tt nsp} must be set with the cellular
structure used by {\tt density} and {\tt interaction}. see
\S\ref{spatial} For a non-cellular model, it contains just the 1
component, equal the {\tt density.size}.


\bibliographystyle{plain}
\bibliography{rus}
\end{document}

